% MAS1004 Assignment 2: long report.
%
% You fill this in with your agent, in your repository. AGENTS.md tells the
% agent how. Red text in the PDF is what still has to be replaced.
%
% - Pictures go in figures/. Copy them from results/, which is not in git.
% - Outputs go in outputs/: copy experiments.csv, LOG.md and the outputs of
%   check_submission.py and curate.py count here. \outputfile{outputs/...}
%   prints a file exactly as it is. Every number in this report has to come
%   from such an output.
%
% To make the PDF: compress this folder into a zip, and on Overleaf choose
% New Project, Upload Project. Set Menu, Compiler to XeLaTeX, and Recompile.
%
% Lines marked % NOAH are Noah's own words, word for word. Everything else is
% written by the agent and only states facts from the outputs.

\documentclass[11pt]{article}
\usepackage{mas1004}

\reporttitle{Assignment 2: Long report}
\studentname{Noah Andreas Jozef Ruijpers}
\studentid{G20260448}

\begin{document}
\makeheader

Leaderboard nickname: stroopwafel123. Final releases: track-a-3 (Track A) and
track-b-1 (Track B).

% ---------------------------------------------------------------------------
\section*{Track A}

All runs below are \texttt{official.py} runs on Colab (250,000 photos through
the network, seed 0 unless stated), with the numbers from
\texttt{experiments.csv}.

\begin{center}\footnotesize
\begin{tabular}{lp{0.6\linewidth}r}
\toprule
tag & what changed from the run before & validation \\
\midrule
baseline & nothing: the starter recipe, run with seeds 0 and 1 & 0.4965 / 0.4918 \\
balanced & \texttt{sample\_weights}: every city drawn equally often & 0.5102 \\
curated  & the balanced recipe, with 51 training photos dropped (28 Busan, 23 Osaka) in \texttt{selection.csv} & 0.5103 \\
\bottomrule
\end{tabular}
\end{center}

Validation per city:

\begin{center}\footnotesize
\begin{tabular}{lrrrrrr}
\toprule
run & average & seoul & busan & tokyo & osaka & taipei \\
\midrule
baseline, seed 0 & 0.4965 & 0.5717 & 0.3734 & 0.6377 & 0.3402 & 0.5595 \\
baseline, seed 1 & 0.4918 & 0.5741 & 0.3827 & 0.6150 & 0.3741 & 0.5132 \\
balanced         & 0.5102 & 0.5306 & 0.4604 & 0.5032 & 0.4591 & 0.5979 \\
curated          & 0.5103 & 0.5425 & 0.4676 & 0.5161 & 0.4361 & 0.5894 \\
\bottomrule
\end{tabular}
\end{center}

What I tried, why, and what I think each result means, as I wrote it in my
lab notebook:

% NOAH: the file below is Noah's Track A entries from LOG.md, exactly as they are.
\outputfile{outputs/log_track_a.txt}

\subsection*{Validation and leaderboard}

My final Track A release is track-a-3: the balanced recipe
(\texttt{track\_a/recipe.py}) with \texttt{track\_a/selection.csv}, the
curated run. The validation scores are from my own runs; the leaderboard
scores are from the leaderboard.

\begin{center}\footnotesize
\begin{tabular}{llll}
\toprule
release & run & validation & leaderboard \\
\midrule
track-a-1 & baseline & 0.4965 (seed 0), 0.4918 (seed 1) & 0.410 \\
track-a-2 & balanced & 0.5102 & 0.416 (range 0.409-0.426) \\
track-a-3 & curated  & 0.5103 & 0.412 (range 0.400-0.433) \\
track-b-1 & b1       & 0.5005 & 0.411 \\
\bottomrule
\end{tabular}
\end{center}

% NOAH
My validation scores are higher than my leaderboard scores. For track-a-2 the grader's own validation score is 0.503, the mean of three seeds, and my leaderboard score is 0.416, so the difference of about 0.09 is far bigger than the seed noise of about 0.005. I think the main reason is that the photos are different. The validation photos come from the same sources as the training photos, and in those sources the frames of one video look much alike, so the model may have seen very similar photos in training. The leaderboard photos are street scenes from YouTube videos whose channels are not in the starter data, with at most 20 frames per video and the same number for every city. So validation is probably too optimistic, and the leaderboard is a harder test.

% ---------------------------------------------------------------------------
\section*{Decisions I made from photos}

I ran \texttt{curate.py suspects} with the balanced model on every fourth
training photo (11,065 photos). For each city it draws the 48 training photos
with the lowest probability for their own city. I looked at the sheets of the
two weakest cities, Busan and Osaka.

\sidebyside{figures/suspects_balanced_seed0_busan_00.png}{\texttt{suspects\_balanced\_seed0\_busan\_00}}%
           {figures/suspects_balanced_seed0_osaka_00.png}{\texttt{suspects\_balanced\_seed0\_osaka\_00}}

\outputfile{outputs/curate_suspects_excerpt.txt}

What I saw in common:

% NOAH
On the Busan and Osaka sheets I saw many photos that show nothing about the city, like food, indoor photos and art. Most of the photos on the sheets were normal street photos, and I kept those.

The rule I wrote before dropping:

% NOAH
\begin{quote}
I drop training photos that show nothing about the city: food, objects, art pieces, indoor photos (homes, malls, restaurants, metro), people, vehicle parts and plain text or logos. I keep normal street photos and buildings, even when they are not special.
\end{quote}

I dropped 28 photos from the Busan sheet and 23 from the Osaka sheet, 51 in
all, with this rule as the reason in \texttt{track\_a/dropped.csv}. The
selection after dropping (\texttt{curate.py count}):

\outputfile{outputs/curate_count.txt}

Scores before and after, with the same recipe and seed:

\begin{center}\footnotesize
\begin{tabular}{lrrrrrr}
\toprule
run & average & seoul & busan & tokyo & osaka & taipei \\
\midrule
balanced (before) & 0.5102 & 0.5306 & 0.4604 & 0.5032 & 0.4591 & 0.5979 \\
curated (after)   & 0.5103 & 0.5425 & 0.4676 & 0.5161 & 0.4361 & 0.5894 \\
\bottomrule
\end{tabular}
\end{center}

The requests I gave my agent for these decisions, word for word:

% NOAH: the file below is Noah's own messages to the agent, word for word.
\outputfile{outputs/agent_requests.txt}

% ---------------------------------------------------------------------------
\section*{Track B}

I made one Track B run, b1: \texttt{train\_b.py} with ResNet18 at
224\,$\times$\,224, starting from torchvision's \texttt{IMAGENET1K\_V1}
weights, 4 epochs, with \texttt{-{}-balance} (every city drawn equally often),
on the starter training photos. It is release track-b-1.

\begin{center}\footnotesize
\begin{tabular}{lrrrrrr}
\toprule
run & validation & seoul & busan & tokyo & osaka & taipei \\
\midrule
b1 & 0.5005 & 0.5027 & 0.4691 & 0.5030 & 0.4403 & 0.5876 \\
\bottomrule
\end{tabular}
\end{center}

% NOAH
I did not collect any extra photos.

My validation score next to my leaderboard score is in the table under
``Validation and leaderboard'' in Track A.

\subsection*{My submission, as the grader sees it}
\outputfile{outputs/check_submission.txt}

% ---------------------------------------------------------------------------
\section*{What made the difference}

% NOAH
In Track A, balanced sampling helped most: +0.0137 over baseline seed 0, more than the seed noise of 0.0047. Dropping 51 photos helped least: 0.5103 against 0.5102, so no measurable effect. In Track B I made only one run (b1), so I cannot say which change helped there.

What I would do next with one more week:

% NOAH: from the "## Next" section of LOG.md, word for word.
With one more week I would look at many more photos of Busan and Osaka, drop the ones that show nothing about the city, and measure that with several seeds. I would also try a different model and a bigger image size in Track B, and add photos of my own.

% ---------------------------------------------------------------------------
\section*{All runs}
\outputfile{outputs/experiments.csv}

\section*{Lab notebook}
\outputfile{outputs/LOG.md}

\end{document}
